% Annexes
\part{Annexes}
\section{Format de fichier INI}
\label{sec:anx_ini}

  Les fichiers sont divisés en sections.

  Chaque section débutte par un nom écrit entre crochets (\texttt{[]}),
  suivi par des lignes du type : \\
  \texttt{option~=~valeur}.

  Les commentaires sont introduits par le caractère dièse (\texttt{\#}).

  Les noms de section et d'option sont généralement \textit{case-sensitive}
  (''sensible à la casse'').

  Ex :
  \begin{lstlisting}
    # Description du format de fichier INI
    [section_1]
    option_1 = exemple

    [EDITOR]
    Prefered = vi
  \end{lstlisting}
